\begin{conclusions}
El presente trabajo de diploma cumplió satisfactoriamente con el objetivo general de 
desarrollar la versión 3.0 de DermaUH, incorporando un nuevo modelo de procesamiento de 
imágenes entrenado con un conjunto de pieles cubanas, mejorando el rendimiento de la 
aplicación y el protocolo de seguridad, e integrando un servicio de correo para la 
gestión de notificaciones. La plataforma resultante constituye una evolución 
significativa respecto a la versión anterior, al abordar las tres limitaciones 
fundamentales identificadas en el estado del arte: la ausencia de un modelo explicativo 
(XAI), el sesgo por entrenamiento con pieles claras y las deficiencias de rendimiento en 
el frontend. A continuación, se exponen los logros alcanzados en relación con cada uno de 
los objetivos específicos planteados.

En cuanto al análisis de herramientas existentes para el diagnóstico dermatológico 
asistido por inteligencia artificial, se examinaron en profundidad tres plataformas 
comerciales (DermEngine, SkinVision y VisualDx) y tres repositorios académicos de 
referencia (ISIC Archive, DERM12345 y HAM10000). El estudio evidenció que ninguna de 
estas soluciones se adapta a las condiciones del sistema de salud cubano, debido a 
barreras económicas (suscripciones anuales entre 50 y 1188 USD), dependencia absoluta de 
conectividad a internet de alta velocidad, entrenamiento mayoritario con imágenes de 
pieles claras (fototipos I-III) y ausencia total de modelos explicativos que permitan al 
especialista comprender el razonamiento detrás de cada diagnóstico. Este análisis 
fundamentó la decisión de desarrollar una solución propia, gratuita, optimizada para 
entornos con conectividad limitada y basada en un modelo XAI alineado con criterios 
clínicos establecidos (regla ABCDE, criterio de Menzies).

Respecto al desarrollo del frontend y backend que supere las deficiencias de la versión 
2.0, se implementó una aplicación web siguiendo el estilo arquitectónico por capas, con 
React en el nivel de presentación (arquitectura basada en características), Django REST 
Framework en el nivel de lógica de negocio (organizado en cuatro servicios independientes:
auth\_service, email\_service, medical\_service e image\_service) y PostgreSQL como sistema 
gestor de base de datos. El frontend redujo significativamente el tiempo de carga, 
superando la latencia de diez segundos reportada en la versión anterior— al eliminar 
animaciones pesadas y optimizar las peticiones asíncronas mediante Axios y hooks 
personalizados. El backend mejoró el protocolo de seguridad mediante autenticación con 
JWT (expiración a 8 horas, token de refresco en cookie HttpOnly) y comunicación cifrada 
con HTTPS (RNF2), además de incorporar un servicio de correo electrónico que gestiona 
notificaciones de registro, activación y restablecimiento de contraseña, cumpliendo así 
con los requisitos RF4 y RF5.

En relación con la mejora de las métricas de los modelos de clasificación mediante 
reentrenamiento con el nuevo conjunto de imágenes de pieles cubanas proporcionado por el 
MINSAP (4576 imágenes distribuidas en cuatro clases), se integró el modelo híbrido 
desarrollado por Daniel Abad, que combina clasificación binaria (melanoma versus no melanoma) 
seguida de clasificación multiclase (carcinoma basal, carcinoma escamoso, otras lesiones). 
Las métricas obtenidas —precisión de 0.788 para melanoma, sensibilidad de 0.820 y F1 de 
0.804, aunque inferiores al $99.7\%$ de sensibilidad reportado en versiones previas 
entrenadas con conjuntos de mayor tamaño (más de 50000 imágenes), representan un resultado 
clínicamente aceptable en el contexto de un dataset limitado y equilibrado mediante 
técnicas de aumentación de datos (rotación a 90 grados). La categoría “otras lesiones” 
presentó métricas más bajas (precisión 0.697, sensibilidad 0.620, F1 0.656), lo que se 
atribuye a la heterogeneidad visual de este grupo, como se refleja en la matriz de 
confusión donde muchas de estas imágenes fueron clasificadas erróneamente como melanoma o 
carcinoma basal.

Finalmente, se logró que la aplicación web clasificara y explicara al especialista la 
decisión del diagnóstico mediante la integración del modelo explicativo (XAI) desarrollado 
por Amanda Noris. Este modelo, basado en SHAP sobre un clasificador XGBoost, extrae las 
características más relevantes de la imagen (asimetría, bordes irregulares, variaciones 
de color, puntos negros, velo azulado) y las presenta al usuario en tres formatos 
complementarios: numérico (lista ordenada de contribuciones), visual (superposición de 
ejes de simetría, paletas de colores, contornos de bordes y máscaras de velo azulado) y 
textual (descripciones generadas alineadas con la regla ABCDE y el criterio de Menzies).
Las pruebas funcionales realizadas (registro de usuario, activación por administrador, 
CRUD completo sobre pacientes, lesiones, diagnósticos e imágenes, compartición de lesiones
entre doctores) demostraron que el sistema cumple con todos los requisitos funcionales 
(RF1 a RF5) y no funcionales (RNF1 a RNF5) definidos en el capítulo 3, quedando pendiente 
la validación en ensayo clínico programada para 2027.

En síntesis, la versión 3.0 de DermaUH constituye una respuesta integral a las brechas 
identificadas en el estado del arte: gratuidad total, entrenamiento con pieles cubanas 
(fototipos IV-VI), integración de un modelo explicativo alineado con la práctica clínica 
y optimización para entornos con conectividad variable. El principal aprendizaje derivado
del proceso de desarrollo es que la precisión de un modelo de clasificación dermatológica 
no puede desligarse de su capacidad para justificar sus predicciones ante el especialista; 
de nada sirve un algoritmo con alta sensibilidad si sus decisiones permanecen opacas. 
Asimismo, se constató que el equilibrio del dataset y la representatividad de las 
categorías menos frecuentes (carcinoma escamoso) son factores críticos que condicionan el 
rendimiento real del modelo, más allá de las métricas globales. Como trabajo futuro, se 
recomienda expandir el conjunto de imágenes mediante colaboraciones con hospitales 
provinciales, implementar un chatbot basado en LLMs (BioGPT con RAG) para consultas 
interactivas, y desplegar el modelo de IA como microservicio independiente para facilitar 
su actualización continua sin afectar al núcleo de la aplicación.
\end{conclusions}
